\scanpage{40}
\uline{Exemples de} $\varphi$ \uline{admissibles}.

Soit $f\in\Gamma(X,\unsure{\mathcal O|\underline O}_X^{\times})$, posons
\[
\varphi_f(P,v)=v(f(P))\qquad\text{si}\quad P\in X(\overline K),\,v\in M_{\overline K}
\]
C'est une fonction \uline{invariante par} $\Gamma$, je dis qu'elle est \uline{admissible}. \deleted{Il suffit de montrer que si $X$ est affine, \deleted{\unsure{comme}} $K[X]$ engendré par $x_1\,-\,x_n$, \inserted{$c\in\operatorname{Div}(M_K)$,} alors $\varphi_f$ est $M_K$-bornée sur $E[X;x_1,-x_n,c]$ i.e. $\exists\,c'\in\operatorname{Div}(M_K)$ tel que}
\[
\deleted{\mathop{\mathrm{Inf}}\limits_i v(x_i(P))\geq c(v)\Longrightarrow |v(f(P)|\leq c'(\unsure{v|w}).}
\]

Par fonctorialité, on est ramené au cas où $X=\mathrm G_{m\,K}$, \deleted{et à vérifier} \inserted{$f$ la section canonique $x$,} \deleted{que} Dans ce cas, \deleted{pour tte} une partie $M_K\unsure{+|-}$bornée de $X(\overline K)\times M_{\overline K}$ \inserted{$=\overline K^*\times M_{\overline K}$} est contenue dans une partie de la forme $E(M,S)$, $S$ partie finie de $M_K$, \deleted{$M\unsure{\geq|>}0$,} $M$ un nb $>0$, formé des $(\unsure{x|P},v)$ tels que
\[
\begin{aligned}
&\text{a) }v(x)\ \deleted{\text{est entier pour }v}\inserted{{}=0\quad\text{i.e. }x\text{ entier}}\quad\text{si}\quad v|K\notin S\\
&\text{b) }\deleted{\unsure{x|v}}\,|v(x)|\leq M\quad\text{si}\quad v|K\in S.
\end{aligned}
\]
% archive sheet number: 39
